\chapter{Literature Review}
Before starting the design of TRIBAL, the team looked at some existing applications that
solve similar problems. This helped the team understand what features are already
available, what is missing, and how TRIBAL can be different from these platforms. This
chapter looks at a few popular apps that are related to roommate finding and activity
matching, and explains their limitations.

\section{Roomster}
Roomster is a roommate and room-finding platform that is mainly used in the United
States and a few other countries. It allows users to post listings for rooms and search for
roommates based on location and budget. However, Roomster does not have any real
compatibility matching system. It mostly works like a classified listing site where users
browse profiles and contact each other manually. There is no scoring system to tell a user
how well they might get along with a potential roommate based on lifestyle habits like
cleanliness, sleep schedule, or smoking


\section{Bumble BFF}
Bumble BFF is a friend-finding mode of the popular Bumble dating app. It uses a swipebased system where users can like or pass on other profiles to find friends with similar
interests. While this works well for general friend matching, it is not designed for
roommate finding. Bumble BFF also does not collect detailed lifestyle information like
budget range or sleep schedule, which are important factors when two people are going to
live together.

\section{Meetup}
Meetup is a platform mainly used to organize and join group activities and events based
on shared interests. It works well for finding activity partners but does not have any
roommate matching feature at all. Meetup is also not very popular in Nepal, and most
local activity groups still rely on Facebook groups or word of mouth instead.

\section{Facebook Groups}
Facebook Groups are widely used in Nepal for both roommate searching and activity
planning. Many people post in groups like “Room for Rent Kathmandu” or “Hiking
Buddies Nepal” to find roommates or activity partners. The biggest problem with this
method is that there is no structured system behind it. Anyone can post anything, there is
no compatibility check, and there is no way to verify whether the listed information is
true. This often leads to mismatched roommates and unsafe meetups with strangers.

\section{Roomi}
Roomi is another roommate-finding application that allows users to filter listings by
location, price, and a few basic preferences. It is more focused on the United States
market and is not commonly used in Nepal. Like Roomster, Roomi does not use any
weighted scoring system, so users still have to manually go through profiles and decide
compatibility on their own.

\section{SpareRoom}
SpareRoom is a roommate listing platform used mostly in the United Kingdom. It allows
detailed filtering of listings, such as gender preference and smoking status, but the
matching is still based on simple filters rather than an actual compatibility score. There is
no algorithm that combines multiple lifestyle factors into a single match percentage.

\section{Summary of Limitations}
After looking at these platforms, the team noticed a few common limitations:
\begin{itemize}
\item Most platforms either focus only on roommate listing or only on activity and
event matching, but not both.
\item None of these apps use a proper compatibility scoring system that takes multiple
lifestyle factors into account at once.
\item Apps that are popular internationally, like Roomster, Roomi, and SpareRoom, are
not commonly used or trusted in Nepal.
\item Local methods like Facebook groups do not provide any safety or compatibility
checks, which can lead to bad roommate experiences.
\item There is no platform that allows users to both find a compatible roommate and
discover local activity partners from one single application.
\end{itemize}
\section{How Tribal is different}
TRIBAL tries to solve these problems by combining roommate matching and activity and
event discovery into a single mobile application. Instead of relying on manual browsing,
TRIBAL uses a weighted compatibility scoring algorithm that compares onboarding
answers from two users and generates a compatibility score out of 100. This score
considers multiple lifestyle factors together, such as cleanliness, budget, sleep schedule,
noise level, smoking, guests, interests, study habit, food, and drinking preference, instead
of just one or two filters.
By keeping the matching system rule-based instead of using machine learning, TRIBAL
is able to give users a transparent and explainable match score. A user can see why their
compatibility score with another person is high or low, which is something most existing
platforms do not provide. This approach is also more practical to build and test within the
time frame of a minor project, while still giving meaningful results to users.